% !TeX root = ../papers/nc-p3.tex
\section{二十次入力の五次対称行列式と全ての隣接辺の二項表示}
\label{sec:quadric-conductor-quintic}
\begin{theorem}\label{thm:quadric-conductor-quintic}
二十次の整な一般被約余因子曲線の入力 $X$ には、二十印で二重に消える五次式がスカラーを除き一つだけある。
これは五行対称一次行列 $A(x)$ の既約行列式であり、階数三以下スキームは $X$ そのものである。
さらに $X$ を一方の射影とする\emph{任意の}有限な六双一次完全交叉辺を取る。
入力座標を固定し、五次行列の五空間を二項四次空間と同定する基底を選べば、
$A(x)$ は二項八次汎関数 $\lambda_x$ の中央Hankel行列となる。
その辺のもう一方の座標を二項三次式 $p$ として、辺の六式は
\[
                   \lambda_x(p f)=0\qquad(f\in S_5)
\]
に等しい。ここで $S_j$ は二項 $j$ 次式の空間である。
\end{theorem}
\begin{proof}[正規化環の対称表示]
$B=\bigoplus_{n\ge0}H^0(Q,\mathcal O_Q(2n))$ を入力の四変数環上の加群とする。
中間コホモロジー消滅によりCohen--Macaulay、標準加群は $B(-1)$ である。
三つの一般入力一次式でNoether正規化 $P_0$ を取り、
$B$ はその上自由で、基底次数は $0,1^6,2$ である（Hilbert関数 $(2n+1)^2$）。
有限双対 $\operatorname{Hom}_{P_0}(B,P_0(-2))\simeq B$ により、
$B$ は $P_0(-2)$ に値をとる完全対称積対合を持つ。
実際 $1$ に対応する関数を積に施せばよい。残る一次式を $t$ とし、その乗法行列を $T_0$ とすると、
$tI-T_0$ はこの対合により対称表示となる。単元成分は対応する非退化定数ブロックを
行列の合同で分離して消去できる（次数零の特殊化で対になる次数を合わせる）。
具体的には $t\cdot1$ が六つの一次基底の非零元であり、双対で次数二基底と一次基底が対をなす
二行二列の単元ブロックを一つ消す。次数差から他の単元はない。
従って最小生成元は $1$ および一次の五元で、対称表示は
\[
           B_0(x)=\begin{pmatrix}c_3(x)&q_2(x)^t\\q_2(x)&A(x)\end{pmatrix}
\]
となる。添字は次数、$A$ は五行対称一次である。
また生成元次数零・一の本数はそれぞれ一・五で一次関係はなく、
自己双対では生成次数 $j$ と関係次数 $3-j$ が対になることからもこの型を確認できる。

$D$ 上では $1$ が有限繊維の環で非零であり、対称行列 $B_0(x)$ の像に第一基底ベクトルは入らない。
従って核に第一成分が非零のベクトルがあり、単元第一成分で消去すれば
$\operatorname{rank}B_0(x)=\operatorname{rank}A(x)$ である。
特に一般の $D$ 上で $A$ は可逆、$\det A$ は非零である。
逆に $A$ の階数が三以下なら $B_0$ の階数は五以下で $x\in D$。
よって $A$ の階数三以下は正規化繊維の長さ三の有限集合で、二十印を含む。
対称五行の余因子複体は定理\ref{thm:double-point-opposite-reconstruction}と同じ
交代行列・跡零行列・対称行列の三写像であり、階数落ちが余次元三なので完全である。
従ってそのイデアルは飽和し分解の自由階数と生成次数は順に
$(15,4),(24,5),(10,6)$。長さは二十で、二十印が既に被約なのでこれで尽きる。
従って各印で階数は三であり、余因子十五個は全四次消滅式の基底になる。

導手 $I_C$ は $D$ 上で正規化の $\mathcal O_Q(-10)$ の押出しであり、$I_C(5)$ には一切断だけある。
$A$ は $C$ 上で階数四以下なのでこの切断は $F=\det A$ である。
従って $F$ の正規化上の零因子は導手因子そのもの。
$F$ に次数四以下の因子があれば、それと八次 $D$ との曲線交叉は台が $C$ に入り、
その因子は $C$ を含んでしまう。定理\ref{thm:cofactor-integral-octic}に反するので $F$ は既約である。
また任意の二重消滅五次式は $5\deg C=100<120$ の重複度評価から $I_C$ に入る。
$F$ は階数三の二十点で二重に消えるので所望の一意性を得る。
\end{proof}

\begin{proof}[隣接辺が与えるPlücker写像]
辺の六行四列一次行列を $E(x)$ とする。十五個の最大小行列式は有限退化の
Eagon--Northcott分解から $I_X(4)$ の基底である。
従って余因子基底との間に線形同型
$\Phi:\operatorname{Sym}^2 k^5\longrightarrow\bigwedge^2 k^6$ があり、
$\Phi(\operatorname{adj}A(x))$ は $E(x)^t$ の核二平面のPlücker点を表す（六空間の体積形式を選ぶ）。
$F$ の一般点の核ベクトルを $p\in\mathbb P^4$ とすれば、$\Phi(p^2)$ は分解可能である。

全 $p$ でそうなることを示す。$K(p)$ を五行四列 $[A_i p]$ とする（$A_i$ は $A$ の係数行列）。
その四次最大小行列式を符号付きで並べて $b(p)$ とする。
$A(x)p=0$ の入射は、既約五次上で一般に同型であり、二十印上だけが直線繊維なので全次元は二。
$\operatorname{rank}K(p)\le r\le3$ の各点では $x$ の繊維次元が $3-r$ 以上なので、この軌跡の次元は $r-1$ 以下である。
最大階数落ちの期待余次元上界二によりそのスキーム $Z$ は純二次元で、
Hilbert--Burchにより飽和イデアルは五個の四次式 $b_i$ で生成される。
また $Z$ は整である。階数三の開集合では入射が $p$ から $x$ を一意に与え、
一般の $x$ は $F$ の階数四の滑らかな点なので、入射はそこで被約である。
二十本の例外直線は成分になれず、純性から結論が従う。

$b:\mathbb P^4\dashrightarrow\mathbb P^4$ は四次の双有理対合である。
四式 $z^t A_i p=0$ の入射において $p\notin Z$ なら $z=b(p)$ 一点で、
階数 $r\le3$ の部分の次元は $(r-1)+(4-r)=3$ 以下である。
$p,z$ の対称性とただ一つの四次元成分から対合性が従う。
四次座標に共通因子はない（基点の余次元二）。

$J(p)=\Phi(p^2)$ を六行交代行列とみると、その四行Pfaffianは $Z$ 上で消える四次式である。
従って余Pfaffian行列は $L b(p)$（$L$ は定数線形写像）である。
一般階数が六なら余Pfaffianをもう一度取ると射影的に $J(p)$ へ戻る。
$z=b(p)$ に置き換え $\Phi^{-1}$ を施すと、$\nu_2(b(z))$ が $z$ の二次式で与えられてしまう。
しかしこの写像の原始座標は共通因子のない四次座標の全二次積で次数八であり不可能である。
一般階数四なら $L b(p)$ は核の二平面を表す分解可能な二形式である。
$b$ は支配的なので $\mathbb P(\operatorname{im} L)$ はGrassmann多様体内の線形空間となる。
その全平面は共通ベクトルを含むか、固定三空間に入る（射影直線同士が二本ずつ交わることを使う）。
前者なら全 $J(p)$ に共通核があり、後者ならその固定三空間上で全 $J(p)$ の交代制限は零である
（三空間内に二次元の核を持つ交代形式は零）。
いずれも $p^2$ が対称二次空間を張り $\Phi$ が全射であることに反する。
従って $\Phi$ は二次Veronese四次元多様体全体を $\operatorname{Gr}(2,6)$ に送る。
\end{proof}

\begin{proof}[二項表示と辺の回復]
普遍二平面束の双対をこの $\mathbb P^4$ に引き戻した $E_0$ は六切断で生成され、
$\det E_0=\mathcal O(2)$、切断の外積は全二次式を張る。一般切断の零点が空なら
$E_0=\mathcal O\oplus\mathcal O(2)$ となり、六切断の外積は高々五次元なので不可能。
零点は滑らかな余次元二の曲面であり、随伴で標準束は $\mathcal O(-3)$。
各成分の一般超平面切断は滑らかな整曲線で標準次数が次数の $-2$ 倍なので、
次数一、すなわち各成分は平面である。二平面は $\mathbb P^4$ で交わるから成分は一つだけ。
そのイデアルの二生成分解から $0\to\mathcal O\to E_0\to I_{\mathbb P^2}(2)\to0$
の局所自由な拡張は $\mathcal O(1)^2$ である。
従って二平面は $p$ に線形な二列で与えられ、任意の非零一次結合が $k^5\to k^6$ の単射である。
一次結合の底 $\mathbb P^1$ 上で $\mathcal O(-1)^5\to\mathcal O^6$ の余核は $\mathcal O(5)$ で、
六切断は五次式の全基底である。その全評価の核を取れば二列は $p\mapsto s p,t p$ の形になる。
従って $\Phi$ は $(sp)\wedge(tp)$ を $p^2$ に戻すBezout線形同型の逆である。

$A(x)$ が可逆な開集合では、その逆が二つの五次式 $f,g$ のBezout行列となる。
この逆行列 $A(x)$ はHankelであり $\langle f,g\rangle$ は対応する
$\lambda_x\in S_8^*$ の $S_5\to S_3^*$ 収縮の核である。
実際 $f$ が五つの単純根 $r_i$ を持つ座標では、逆Bezout行列は
評価ベクトルの平方の和で、係数は共通スカラーを除き $1/(f'(r_i)g(r_i))$ である。
$f$ の五根でBezout対合を対角化すればこの式を得る。
$g$ を掛けた三次式の評価和もLagrange補間により零である。
Hankel条件とこの二つの収縮消滅は多項式恒等式なので全ての可逆なBezout行列に成立する。
従って $A(x)$ の全線形系がHankelであり、$\lambda_x$ は $x$ に線形である。
$S_3\to S_5^*$ の二項収縮行列の像は一般に $E(x)$ の像と同じである。
ここでこの収縮が階数三以下なら、消される非零三次式 $p$ によって $sp,tp$ が
中央行列の核に入り $\operatorname{rank}A(x)\le3$、従って $x\in X$。
両行列の像は余次元一で同じ部分束であり、いずれも $\mathcal O(-1)^4$ と同型な反射層なので
全体で一致する（共通分数空間内で余次元一の格子の交わりを取る）。
よって四列の定数基底変更で両者は一致し、全ての隣接辺が主張の形となる。
\end{proof}
\path{computations/verify_quadric_conductor_quintic.py} は有理錐射影で五個の二次関係から対称一次ブロックを復元し、
その十五個の四次余因子と有限辺の最大小行列式の包を良い素数で比較する。
